\chapter{Bevezetés}

A videójáték fejlesztés az elmúlt években több millárd dolláros iparággá nőtte ki magát. Nem csak a rohamosan fejlődő hardverestechnológiának, de a szoftveres ágon tett előrelepéseknek is köszönhetően. A 2000-es évek elején a videójáték fejlesztés mindössze a programozók (matematikusok) számára egy hobbi volt, ami akkoriban a hardver határainak a feszegetéséről szólt. Manapság a játékmotoroknak (game engine) és játék könyvtáraknak (game library) hála bárki elkezdhet játékokat fejleszteni, viszont ez nem azt jelenti, hogy sokkal egyszerűbb lenne a fejlesztők dolga. A folyamatosan megújuló ipar és a játékfejlesztés elérhetőségének a növekedése az elvárásokat is megnö\-velte a játékok iránt.\\

Mivel a számítástechnika és a matematika rokon tudományok, ezért sok matematikai modellt alkalmazunk a programozás során. Egyike ezeknek a modelleknek az ún. "állapotgépek". Az állapotgépeket bővebben kifejtem majd a \hyperref[chap:FSM]{\ref{chap:FSM}. fejezet}-ben. Azon esetekben ajánlatos az alkalmazásuk, ahol fontos a könnyű bővíthetőség, átláthatóság és a programozott probléma szétszedhető esetekre, állapotokra. A legnépszerűbb példa állapotgépek hasz\-nálatára az a videójáték fejlesztéshez kapcsolódik. A kezdetekben ez nem volt olyan fontos, mivel maguk a játékok egyszerűbbek voltak, lásd: Doom (id Software 1993), Wolfenstein (id Software 1992), Elite (Acornsoft 1984) ahol a legnagyobb kihívást a grafika megoldása és annak helyes kirajzolása jelentette. Manapság, ahol ún. "Omni-Movement" van a játékokban, ami, ahogyan Charles \cite{charles2024} is megfogalmazta a posztjában, egy olyan mozgásfajta, ahol a játékos teste a fejétől függetlenül mozog, tehát a test és a fej (kamera) mutathat két különböző irányba, valamint a felhasználó képes számos mozdulatot végrehajtani a karakterrel és minden mozdulat helyzetfüggő. Ezen esetben a játékosoknak közelről kell figyelni a mozgásukat és cselekedeteiket, és jól definiált struktúrákat kell létrehozni egyes mozgássorozatoknak. A játékosok képesek különböző mozdulatok végre\-hajtására, többek között: futás, ugrás, csúszás, vetődés, stb. Ezek mind precíz inputokat követelnek meg a felhasználótól. Ennek a megvalósításához az állapotgépek használata elengedhetetlen, mind a mozgásrendszer megtervezéséhez, mind az animációk helyes lejátszásához.\\

A témámat, viszont nem az Omni-Movement inspirálta, az én projektemhez csak példaképpen kapcsolódik, hogy demonstráljam egy sokkal bonyolultabb koncepción is a témámat és bizonyítsam, hogy az iparban is használa\-tos. Sokkal inkább a Shibuya-punk esztétika \cite{Yoko2023} királya és a mai napig a Sega egyik legkevésbé elismert kiadása a: Jet Set Radio (Sega 2000), és az az inspirálta Bomb Rush Cyberfunk (Team Reptile 2023). Ezekről sokan nem hallottak, viszont, ha a "Tony Hawk Pro Skater" sorozatot említem, akkor az már több embernek fog ismerősen hangzani. A koncepció ugyan az mind a kettőnél: gurulj és csinálj menő trükköket (miközben próbálod magad nem összetörni, természetesen). Ha a felszínt nézzük is már eléggé bonyolultnak néz ki a helyzet, mind mozgás, mind pontszám számítás szempontjából; és elkezd  járni az agyunk azon, hogy mégis mennyi állapotot kellet nyomon követniük a fejlesztőknek, hogy ezeket megvalósítsák. Annak érdekében, hogy egy kicsit mélyebbre ássak és megválaszoljam ezt a kérdést úgy döntöttem, hogy magam is nekiállok és létrehozok egy mozgásrendszert, amiben tesztelhetjük, hogy mennyire is érné meg állapotgépeket használni, miért és mennyivel eredményezne szebb, jobban strukturált, átláthatóbb, könnyebben bővíthető kódot, mint egy hagyományos if-else ágakból álló megol\-dás.\\

A szakdolgozat végére szeretnék egy teljesen működőképes és játszható mozgásrendszer-prototípust elkészíteni, amely bemutatja az állapotgépek fon\-tosságát a mozgásrendszerek kialakításában és a játékfejlesztés más terein. Mind ezt a Shibuya-punk esztétikában többnyire magam által elkészített modellekből.\\